\begin{thebibliography}{99}
\small

\bibitem{Abramenko}
P.~Abramenko,
Group actions on twin buildings,
\emph{Bulletin of the Belgian Mathematical Society -- Simon Stevin}
\textbf{3} (1996), no.~4, 391--406.
\href{https://eudml.org/doc/119921}{EuDML:119921}.

\bibitem{AbramenkoBrown}
P.~Abramenko and K.~S.~Brown,
\emph{Buildings: Theory and Applications},
Graduate Texts in Mathematics \textbf{248}, Springer, 2008.
\href{https://doi.org/10.1007/978-0-387-78835-7}{doi:10.1007/978-0-387-78835-7}.

\bibitem{AVM2001}
P.~Abramenko and H.~Van Maldeghem,
1-twinnings of buildings,
\emph{Mathematische Zeitschrift} \textbf{238} (2001), no.~1, 187--203.
\href{https://doi.org/10.1007/PL00004899}{doi:10.1007/PL00004899}.

\bibitem{BH1999}
M.~R.~Bridson and A.~Haefliger,
\emph{Metric Spaces of Non-Positive Curvature},
Grundlehren der mathematischen Wissenschaften \textbf{319}, Springer, 1999.
\href{https://doi.org/10.1007/978-3-662-12494-9}{doi:10.1007/978-3-662-12494-9}.

\bibitem{CL2025}
M.~Chapman and A.~Lubotzky,
Stability of homomorphisms, coverings and cocycles II: Examples,
applications and open problems,
\emph{Advances in Mathematics} \textbf{463} (2025), article 110117.
\href{https://doi.org/10.1016/j.aim.2025.110117}{doi:10.1016/j.aim.2025.110117}.

\bibitem{DD2024}
Y.~Dikstein and I.~Dinur,
Coboundary and cosystolic expansion without dependence on dimension or degree,
in \emph{APPROX/RANDOM 2024}, LIPIcs \textbf{317}, 62:1--62:24.
\href{https://doi.org/10.4230/LIPIcs.APPROX/RANDOM.2024.62}{doi:10.4230/LIPIcs.\allowbreak APPROX/RANDOM.2024.62}.

\bibitem{EK2024}
S.~Evra and T.~Kaufman,
Bounded degree cosystolic expanders of every dimension,
\emph{Journal of the American Mathematical Society}
\textbf{37} (2024), no.~1, 39--68.
\href{https://doi.org/10.1090/jams/1019}{doi:10.1090/jams/1019}.

\bibitem{GdPVB2025}
L.~Grave de Peralta and I.~Valentiner-Branth,
High-dimensional expanders from Kac--Moody--Steinberg groups,
\emph{European Journal of Combinatorics} \textbf{126} (2025), article 104131.
\href{https://doi.org/10.1016/j.ejc.2025.104131}{doi:10.1016/j.ejc.2025.104131}.

\bibitem{Gromov2010}
M.~Gromov,
Singularities, expanders and topology of maps. Part 2: From combinatorics
to topology via algebraic isoperimetry,
\emph{Geometric and Functional Analysis} \textbf{20} (2010), no.~2, 416--526.
\href{https://doi.org/10.1007/s00039-010-0073-8}{doi:10.1007/s00039-010-0073-8}.

\bibitem{Iwasa}
R.~Iwasa,
Homology pro stability for Tor-unital pro rings,
\emph{Homology, Homotopy and Applications} \textbf{22} (2020), no.~1, 343--374.
\href{https://doi.org/10.4310/HHA.2020.v22.n1.a20}{doi:10.4310/HHA.2020.v22.n1.a20}.

\bibitem{KKL2016}
T.~Kaufman, D.~Kazhdan, and A.~Lubotzky,
Isoperimetric inequalities for Ramanujan complexes and topological expanders,
\emph{Geometric and Functional Analysis} \textbf{26} (2016), no.~1, 250--287.
\href{https://doi.org/10.1007/s00039-016-0362-y}{doi:10.1007/s00039-016-0362-y}.

\bibitem{KM2021}
T.~Kaufman and D.~Mass,
Unique-neighbor-like expansion and group-independent cosystolic expansion,
in \emph{ISAAC 2021}, LIPIcs \textbf{212}, 56:1--56:17.
\href{https://doi.org/10.4230/LIPIcs.ISAAC.2021.56}{doi:10.4230/LIPIcs.ISAAC.2021.56}.

\bibitem{KO2021}
T.~Kaufman and I.~Oppenheim,
Coboundary and cosystolic expansion from strong symmetry,
in \emph{ICALP 2021}, LIPIcs \textbf{198}, 84:1--84:16.
\href{https://doi.org/10.4230/LIPIcs.ICALP.2021.84}{doi:10.4230/LIPIcs.ICALP.2021.84}.

\bibitem{KO2023}
T.~Kaufman and I.~Oppenheim,
High dimensional expanders and coset geometries,
\emph{European Journal of Combinatorics} \textbf{111} (2023), article 103696.
\href{https://doi.org/10.1016/j.ejc.2023.103696}{doi:10.1016/j.ejc.2023.103696}.

\bibitem{KOW2025}
T.~Kaufman, I.~Oppenheim, and S.~Weinberger,
\emph{Coboundary expansion of coset complexes}, 2024;
conference version in \emph{STOC 2025}.
\href{https://arxiv.org/abs/2411.02819v1}{arXiv:2411.02819v1}.

\bibitem{LM2006}
N.~Linial and R.~Meshulam,
Homological connectivity of random 2-complexes,
\emph{Combinatorica} \textbf{26} (2006), no.~4, 475--487.
\href{https://doi.org/10.1007/s00493-006-0027-9}{doi:10.1007/s00493-006-0027-9}.

\bibitem{LLR2019}
A.~Lubotzky, Z.~Luria, and R.~Rosenthal,
Random Steiner systems and bounded degree coboundary expanders of every dimension,
\emph{Discrete \& Computational Geometry} \textbf{62} (2019), no.~4, 813--831.
\href{https://doi.org/10.1007/s00454-018-9991-2}{doi:10.1007/s00454-018-9991-2}.

\bibitem{LM2015}
A.~Lubotzky and R.~Meshulam,
Random Latin squares and 2-dimensional expanders,
\emph{Advances in Mathematics} \textbf{272} (2015), 743--760.
\href{https://doi.org/10.1016/j.aim.2014.12.015}{doi:10.1016/j.aim.2014.12.015}.

\bibitem{LMM2016}
A.~Lubotzky, R.~Meshulam, and S.~Mozes,
Expansion of building-like complexes,
\emph{Groups, Geometry, and Dynamics} \textbf{10} (2016), no.~1, 155--175.
\href{https://doi.org/10.4171/GGD/346}{doi:10.4171/GGD/346}.

\bibitem{MW2009}
R.~Meshulam and N.~Wallach,
Homological connectivity of random $k$-dimensional complexes,
\emph{Random Structures \& Algorithms} \textbf{34} (2009), no.~3, 408--417.
\href{https://doi.org/10.1002/rsa.20238}{doi:10.1002/rsa.20238}.

\bibitem{Morrow}
M.~Morrow,
Pro unitality and pro excision in algebraic $K$-theory and cyclic homology,
\emph{Journal f\"ur die reine und angewandte Mathematik} \textbf{736} (2018), 95--139.
\href{https://arxiv.org/abs/1404.4179v2}{arXiv:1404.4179v2}.

\bibitem{OP2022}
R.~O'Donnell and K.~Pratt,
High-dimensional expanders from Chevalley groups,
in \emph{CCC 2022}, LIPIcs \textbf{234}, 18:1--18:26.
\href{https://doi.org/10.4230/LIPIcs.CCC.2022.18}{doi:10.4230/LIPIcs.CCC.2022.18}.

\bibitem{OVB}
I.~Oppenheim and I.~Valentiner-Branth,
\emph{New cosystolic high-dimensional expanders from KMS groups},
2025, revised September 9, 2026; to appear in \emph{Mathematische Annalen}.
\href{https://arxiv.org/abs/2504.05823v2}{arXiv:2504.05823v2}.

\bibitem{OVBNonAbelian2025}
I.~Oppenheim and I.~Valentiner-Branth,
\emph{Non-Abelian expansion of congruence KMS complexes}, 2025.
\href{https://arxiv.org/abs/2510.18351v1}{arXiv:2510.18351v1}.

\bibitem{Vas1969}
L.~N.~Vaserstein,
On the stabilization of the general linear group over a ring,
\emph{Mathematics of the USSR-Sbornik} \textbf{8} (1969), no.~3, 383--400.
\href{https://doi.org/10.1070/SM1969v008n03ABEH001279}{doi:10.1070/SM1969v008n03ABEH001279}.

\end{thebibliography}
